\section{Folha de registro da atividade As Chaves do Cofre}
\label{apendice:chaves-cofre}

A folha de registro apresentada a seguir foi elaborada para coletar
informações sobre o processo de resolução da atividade
\textit{As Chaves do Cofre}.

\section{Folha de registro da atividade Água para Todos}
\label{apendice:agua-todos}

A folha de registro apresentada a seguir foi elaborada para coletar
informações sobre o processo de resolução da atividade
\textit{Água para Todos}.

\section{Prompt utilizado na estruturação dos registros com Gemini}
\label{apendice:prompt-gemini}

O processamento dos registros utiliza instruções distintas conforme a estrutura de cada atividade. Além das imagens digitalizadas, o Gemini recebe o texto previamente extraído pelo Google Document AI e uma
descrição dos campos presentes na atividade.

Na atividade \textit{As Chaves do Cofre}, o processamento utiliza uma única imagem, correspondente à folha preenchida pelo estudante. Na atividade \textit{Água para Todos}, são utilizadas duas imagens: a folha da questão, que pode conter marcações realizadas diretamente
sobre o mapa, e a folha de registro.

A resposta do modelo é solicitada em formato JSON e segue uma estrutura previamente definida pelo programa para cada atividade. Os prompts utilizados são apresentados nas subseções seguintes.


\subsection*{Prompt para a atividade As Chaves do Cofre}

Na atividade \textit{As Chaves do Cofre}, o Gemini recebe a imagem JPG da folha preenchida pelo estudante, o texto extraído pelo OCR e a descrição da estrutura da atividade.

O prompt utilizado é apresentado a seguir. Os elementos entre chaves são substituídos automaticamente pelo programa durante o processamento.

\begin{lstlisting}[style=prompt]
Você está analisando uma folha de registro de um estudante
na atividade Bebras:

"{configuracao['nome']}"

A imagem contém elementos impressos da atividade
e respostas produzidas pelo estudante.

As respostas podem conter:

- texto manuscrito;
- alternativas marcadas;
- desenhos;
- setas;
- formas geométricas;
- sequências;
- riscos;
- outras representações.

O texto abaixo foi extraído pelo Google Document AI.

--- INÍCIO DO OCR ---

{texto_ocr}

--- FIM DO OCR ---

ESTRUTURA DA ATIVIDADE:

{configuracao["descricao"]}

REGRAS:

- Use a imagem como principal fonte.
- Use o OCR somente como apoio.
- Diferencie elementos impressos das respostas do estudante.
- Não invente informações.
- Não determine a resposta correta da atividade.
- Não avalie o desempenho do estudante.
- Não atribua habilidades de Pensamento Computacional.

IMPORTANTE SOBRE DESENHOS E MARCAÇÕES:

- Descreva apenas aquilo que é diretamente visível.
- Não atribua intenção ao estudante.
- Não explique por que o estudante fez determinada marcação,
  a menos que isso esteja explicitamente escrito por ele.
- Não transforme uma marcação visual em uma conclusão
  sobre o raciocínio do estudante.
- Não use expressões como "para verificar",
  "com o objetivo de", "indicando que" ou semelhantes
  para explicar a intenção do estudante,
  a menos que essas palavras tenham sido escritas
  pelo próprio estudante.

Exemplo:

CORRETO:

"Ponto 4 circulado. Há linhas traçadas entre o ponto 4
e outros nós."

INCORRETO:

"O estudante circulou o ponto 4 para verificar se todas
as vilas estavam a duas estradas."

STATUS DAS QUESTÕES:

Use "ok" SOMENTE quando a informação estiver
diretamente legível ou visualmente identificável
sem necessidade de interpretação incerta.

Use "revisar" quando:

- alguma parte estiver ambígua;
- existir mais de uma interpretação possível;
- uma forma ou símbolo não puder ser identificado
  com total segurança;
- apenas parte de uma sequência estiver clara;
- você sentir necessidade de usar palavras como
  "parece", "possivelmente", "provavelmente",
  "aparentemente" ou "talvez".

Se houver dúvida, prefira "revisar" em vez de "ok".

Use "ilegivel" quando houver conteúdo produzido
pelo estudante, mas não for possível interpretá-lo.

Uma resposta realmente em branco deve ser null.
Uma resposta em branco pode ter status "ok",
pois foi possível verificar que o campo está vazio.

NOME DO ESTUDANTE:

- Se estiver claramente legível:
  nome_status = "ok".

- Se houver uma leitura provável, mas com dúvida:
  registre a leitura provável e use
  nome_status = "revisar".

- Se houver um nome escrito mas não for possível lê-lo:
  nome_estudante = null e
  nome_status = "ilegivel".

- Se o campo estiver vazio:
  nome_estudante = null e
  nome_status = "ausente".

Preserve o sentido original das respostas.

Corrija erros do OCR somente quando a imagem
permitir confirmar claramente o conteúdo.

Retorne somente JSON.
\end{lstlisting}


\subsection*{Estrutura informada para As Chaves do Cofre}

Além do prompt, o programa fornece ao modelo uma descrição dos campos presentes na folha de registro:

\begin{lstlisting}[
    breaklines=true,
    basicstyle=\ttfamily\scriptsize,
    columns=fullflexible,
    keepspaces=true
]
1. Qual alternativa você escolheu?
Opções: A, B, C ou D.

2. Registre a sequência de objetos utilizada.
Pode conter desenhos, formas geométricas e setas.

3. Explique como você encontrou esse caminho.
Pode conter texto, desenhos, setas ou estar em branco.

4. Você testou algum caminho que não funcionou?
Opções: Sim ou Não.
Pode existir uma representação desenhada.

5. Como você verificou que esse era o menor número
possível de gavetas?
Resposta aberta.

6. Qual foi a parte mais difícil da atividade?
Resposta aberta.

7. Depois de conversar com um colega, você mudou
sua resposta ou estratégia?
Opções: Sim ou Não.
Pode conter justificativa.
\end{lstlisting}


\subsection*{Prompt para a atividade Água para Todos}

Na atividade \textit{Água para Todos}, o Gemini recebe duas imagens do mesmo estudante. A primeira corresponde à folha da questão e do mapa, enquanto a segunda corresponde à folha de registro. O texto extraído pelo Google Document AI é obtido da folha de registro e utilizado como apoio à leitura das imagens.

O prompt utilizado é apresentado a seguir.

\begin{lstlisting}[
    breaklines=true,
    basicstyle=\ttfamily\scriptsize,
    columns=fullflexible,
    keepspaces=true
]
Você está analisando o registro de um estudante
na atividade Bebras:

"{configuracao['nome']}"

Você recebeu DUAS imagens do mesmo estudante.

IMAGEM 1:

É a folha da questão/mapa.

Ela contém elementos originalmente impressos,
mas também pode conter anotações feitas pelo estudante,
como:

- caminhos;
- setas;
- riscos;
- círculos;
- números;
- símbolos;
- marcações;
- pequenas anotações.

Essas anotações fazem parte do processo
de resolução do estudante.

IMAGEM 2:

É a folha de registro/resposta.

O texto abaixo foi extraído pelo
Google Document AI da folha de resposta.

REGRAS:

- Use a imagem como principal fonte.
- Use o OCR somente como apoio.
- Diferencie elementos impressos das respostas do estudante.
- Não invente informações.
- Não determine a resposta correta da atividade.
- Não avalie o desempenho do estudante.
- Não atribua habilidades de Pensamento Computacional.

IMPORTANTE SOBRE DESENHOS E MARCAÇÕES:

- Descreva apenas aquilo que é diretamente visível.
- Não atribua intenção ao estudante.
- Não explique por que o estudante fez determinada marcação,
  a menos que isso esteja explicitamente escrito por ele.
- Não transforme uma marcação visual em uma conclusão
  sobre o raciocínio do estudante.
- Não use expressões como "para verificar",
  "com o objetivo de", "indicando que" ou semelhantes
  para explicar a intenção do estudante,
  a menos que essas palavras tenham sido escritas
  pelo próprio estudante.

Exemplo:

CORRETO:

"Ponto 4 circulado. Há linhas traçadas entre o ponto 4
e outros nós."

INCORRETO:

"O estudante circulou o ponto 4 para verificar se todas
as vilas estavam a duas estradas."

STATUS DAS QUESTÕES:

Use "ok" SOMENTE quando a informação estiver
diretamente legível ou visualmente identificável
sem necessidade de interpretação incerta.

Use "revisar" quando:

- alguma parte estiver ambígua;
- existir mais de uma interpretação possível;
- uma forma ou símbolo não puder ser identificado
  com total segurança;
- apenas parte de uma sequência estiver clara;
- você sentir necessidade de usar palavras como
  "parece", "possivelmente", "provavelmente",
  "aparentemente" ou "talvez".

Se houver dúvida, prefira "revisar" em vez de "ok".

Use "ilegivel" quando houver conteúdo produzido
pelo estudante, mas não for possível interpretá-lo.

Uma resposta realmente em branco deve ser null.
Uma resposta em branco pode ter status "ok",
pois foi possível verificar que o campo está vazio.

NOME DO ESTUDANTE:

- Se estiver claramente legível:
  nome_status = "ok".

- Se houver uma leitura provável, mas com dúvida:
  registre a leitura provável e use
  nome_status = "revisar".

- Se houver um nome escrito mas não for possível lê-lo:
  nome_estudante = null e
  nome_status = "ilegivel".

- Se o campo estiver vazio:
  nome_estudante = null e
  nome_status = "ausente".

Preserve o sentido original das respostas.

Corrija erros do OCR somente quando a imagem
permitir confirmar claramente o conteúdo.

Retorne somente JSON.

--- INÍCIO DO OCR ---

{texto_ocr}

--- FIM DO OCR ---

ESTRUTURA DA ATIVIDADE:

{configuracao["descricao"]}

REGRAS:

- Analise as duas imagens em conjunto.
- Use as imagens como principal fonte.
- Use o OCR somente como apoio textual.
- Considere as anotações feitas na folha da questão.
- Diferencie elementos impressos dos elementos
  acrescentados pelo estudante.
- Não invente informações.
- Não determine a resposta correta.
- Não avalie o estudante.
- Não atribua habilidades de Pensamento Computacional.
- Apenas transcreva e descreva aquilo que foi registrado.
- Preserve o sentido original das respostas.
- Uma resposta realmente em branco deve ser null.
- Não confunda resposta em branco com conteúdo ilegível.
- Se houver ambiguidade, use status "revisar".
- Se houver conteúdo ilegível, use status "ilegivel".
- Se estiver claro, use status "ok".
- Se um campo não puder ser determinado com segurança,
  retorne null e continue preenchendo os demais.

Retorne somente JSON.
\end{lstlisting}


\subsection*{Estrutura informada para Água para Todos}

A descrição da folha de registro fornecida ao modelo é a seguinte:

\begin{lstlisting}[
    breaklines=true,
    basicstyle=\ttfamily\scriptsize,
    columns=fullflexible,
    keepspaces=true
]
1. Qual alternativa você escolheu?
Opções: A, B, C ou D.

2. Mostre como você analisou os caminhos do mapa
para chegar à resposta.
Pode conter desenhos, marcações, setas, texto ou símbolos.

3. Quais pontos você chegou a considerar antes
de escolher a resposta final?
Podem ser marcados os pontos 1, 2, 3 e 4.
Também pode ser marcada a opção indicando que
apenas a resposta final foi considerada.

4. Você testou algum ponto que não funcionou?
Opções: Sim ou Não.
Pode haver explicação ou representação do ponto testado.

5. Como você verificou que, com o ponto escolhido,
todas as vilas ficam a no máximo duas estradas
de um centro de distribuição?
Pode conter texto, desenhos ou marcações.

6. Qual foi a parte mais difícil da atividade?
Resposta aberta.

7. Depois de conversar com um colega, você mudou
sua resposta ou estratégia?
Opções: Sim ou Não.
Pode conter justificativa.
\end{lstlisting}


\subsection*{Estrutura da saída}

A saída do Gemini é restringida por um esquema JSON definido
previamente para cada atividade. O esquema determina os campos que
podem ser retornados e seus respectivos tipos de dados.

Em ambas as atividades, a estrutura contém a identificação da
atividade, o nome do estudante, o estado de leitura do nome, o ano e
os campos correspondentes às sete questões da folha de registro.

O campo \texttt{nome\_status} pode assumir os valores
\texttt{ok}, \texttt{revisar}, \texttt{ilegivel} ou
\texttt{ausente}. Para as questões, são utilizados os estados
\texttt{ok}, \texttt{revisar} e \texttt{ilegivel}.

O valor \texttt{null} pode ser utilizado quando um campo está em branco
ou quando determinada informação não pode ser preenchida com segurança,
conforme as regras estabelecidas no prompt. Os campos específicos
dependem do tipo de questão e podem armazenar alternativas, textos,
listas, valores booleanos ou descrições de representações produzidas
pelo estudante.

O programa solicita ao modelo uma resposta com o tipo
\texttt{application/json} e utiliza o esquema correspondente à
atividade para restringir a estrutura retornada. Após o processamento,
o JSON é normalizado e seus dados são utilizados para gerar a planilha
empregada na etapa de conferência.